%%% conclusion.tex — Заключение
\section*{ЗАКЛЮЧЕНИЕ}
\addcontentsline{toc}{section}{ЗАКЛЮЧЕНИЕ}

Проведённое исследование подтвердило, что сетевая стеганография представляет реальную и труднообнаруживаемую угрозу для информационной безопасности корпоративных сетей. Анализ деятельности APT-группировок (Turla, Equation Group, APT29) показал, что скрытые каналы в протоколах TCP/IP активно эксплуатируются для управления вредоносным ПО и эксфильтрации данных, при этом существующие средства защиты (DPI, IDS/IPS, DLP) не обеспечивают их надёжного обнаружения.

Сравнительный анализ сетевой и контейнерной стеганографии выявил, что встраивание данных в поля заголовков пакетов, несмотря на существенно меньшую ёмкость по сравнению с мультимедийными контейнерами, обладает принципиальным преимуществом --- оно не создаёт новых информационных потоков и не требует передачи дополнительных файлов, что делает его практически невидимым для стандартных средств мониторинга. Данное свойство определяет целесообразность сетевой стеганографии для передачи малых объёмов критически важных данных в контролируемых сетевых средах.

Разработанная система \texttt{NetStego} реализует пять стеганографических каналов с ёмкостью от 2 до 1470~байт на пакет, покрывая диапазон от высокоскрытных (IP~ID, TCP~ISN, TCP~Timestamp) до высокоёмких (ICMP Payload) каналов. Эффективная пропускная способность варьируется от 20~байт/с до 115{,}9~Кбит/с при скорости 10~PPS, что подтверждает фундаментальный компромисс между скоростью передачи и скрытностью канала.

Применение криптографического алгоритма AES-256-GCM с деривацией ключа Argon2id и двухуровневого контроля целостности (CRC32 на уровне чанков, SHA-256 на уровне файла) обеспечивает конфиденциальность, аутентичность и целостность передаваемых данных. Протокол фрагментации NST с поддержкой сборки чанков не по порядку доставки обеспечивает устойчивость к переупорядочиванию пакетов в сети.

Модуль обнаружения скрытых каналов объединяет четыре метода детекции: статистический анализ (9~тестов, включая канал-специфичные проверки), алгоритм Isolation Forest (unsupervised, 9~признаков с учётом энтропии payload), алгоритм Random Forest (supervised, обучение на размеченных данных) и сигнатурный поиск. Результаты тестирования (269~тестов) продемонстрировали, что комбинация методов позволяет обнаруживать все реализованные каналы с F1-мерой от 0{,}80 до 0{,}90+ при уровне ложноположительных срабатываний не более 10\,\%.

Модуль анализа содержимого payload реализует обнаружение вредоносных вложений, скрытых в стеганографическом трафике. Сканирование извлечённых данных на наличие 20~сигнатур в 7~категориях (исполняемые файлы, скрипты, шеллкод, архивы, макро-документы, веб-атаки, обфускация) позволяет не только выявить факт использования скрытого канала, но и классифицировать тип скрытых данных для оценки уровня угрозы.

Реализованный механизм избыточности (дублирование чанков с дедупликацией по порядковому номеру) обеспечивает устойчивость к потере пакетов для высокоёмких каналов: при трёхкратной избыточности данные восстанавливаются при потере до 20\,\% пакетов. Система сбора статистики, интегрированная в конвейеры отправки и приёма, обеспечивает протоколирование каждой сессии в базе данных SQLite с поддержкой экспорта в JSON, CSV и HTML. Модуль визуализации автоматически генерирует 6~типов графиков для демонстрации результатов производительности и точности обнаружения.

Вместе с тем выявлены ограничения системы: отсутствие механизма ARQ для повторной передачи потерянных фрагментов, обнаруживаемость сигнатуры NST в полезной нагрузке высокоёмких каналов, а также уязвимость к активным нарушителям, способным модифицировать подозрительные пакеты. Устранение данных ограничений --- применение помехоустойчивого кодирования (Reed--Solomon), обфускация сигнатур и адаптивная маскировка временных характеристик --- определяет направления дальнейших исследований.

Практическая значимость работы состоит в создании инструмента, пригодного для исследования уязвимостей сетевых протоколов, обучения специалистов по кибербезопасности и тестирования корпоративных систем защиты. Реализованный модуль обнаружения, объединяющий статистический, ML и сигнатурный анализ с модулем обнаружения вредоносных вложений, может быть интегрирован с системами SIEM для автоматического выявления стеганографической активности и классификации скрытых угроз в реальном времени.

\textbf{Перспективы дальнейшего развития:}
\begin{enumerate}[label=\arabic*)]
    \item расширение набора каналов за счёт IPv6 Extension Headers, протоколов HTTP/2 и QUIC;
    \item реализация ARQ-протокола с использованием NACK-флага для повторной передачи;
    \item обфускация сигнатуры NST с помощью XOR-маски, производной от ключа;
    \item замена простого дублирования на коды коррекции ошибок (Reed--Solomon);
    \item использование глубокого обучения (CNN/LSTM) для анализа последовательностей пакетов;
    \item внедрение модуля обнаружения в режиме реального времени на основе высокопроизводительных фреймворков PF\_RING или DPDK;
    \item интеграция с системами SIEM для автоматического оповещения при обнаружении стеганографической активности.
\end{enumerate}
